\chapter{Análisis de Pequeña Señal y Parámetros del Amplificador}
\label{chap:pequena_senal}

En este capítulo se analiza el comportamiento dinámico linealizado para señales de pequeña amplitud mediante el
modelo híbrido equivalente del transistor BJT, deduciendo las expresiones teóricas para la impedancia de entrada
($Z_i$), impedancia de salida ($Z_o$), ganancia de tensión ($A_v$) y ganancia de corriente ($A_i$). Asimismo, se
formalizan los métodos de medición experimental implementados en el laboratorio.

\section{Modelo Híbrido Equivalente del Transistor}

En régimen de pequeña señal, el transistor BJT polarizado en zona activa se modela linealmente despreciando la
conductancia de salida ($h_{oe} \approx 0$) y la realimentación inversa ($h_{re} \approx 0$). En la configuración
colector común, el terminal de colector actúa como masa de señal de corriente alterna.

En la figura~\ref{crkt:pequena_senal} se ilustra el circuito equivalente del amplificador completo.

\begin{figure}[!ht]
  \centering
  \begin{circuitikz}[scale=0.9, transform shape]
  % Terminal de entrada
  \draw (0,0) node[ground] {}
    to[open, v^>=$v_i$] (0,2.5)
    to[short, o-*] (1.5,2.5) coordinate (in_node);

  % Resistencia RB = R1 // R2 a masa de CA
  \draw (in_node) to[R, l=$R_B$, -] (1.5,0) node[ground] {};

  % Modelo híbrido: resistencia de entrada hie entre base y emisor
  \draw (in_node) to[short] (3.0,2.5)
    to[R, l=$h_{ie}$, i=$i_b$] (3.0,1.2) coordinate (emit_node);

  % Fuente controlada hfe*ib inyectando corriente al nodo de emisor
  \draw (5.5,0) node[ground] {}
    to[cI, l=$h_{fe} i_b$] (5.5,2.5)
    to[short] (5.5,1.2)
    to[short, -*] (emit_node);

  % Nodo de emisor conectado a RE y RL en paralelo a masa
  \draw (emit_node) to[short] (7.0,1.2) coordinate (aux_e);
  \draw (aux_e) to[R, l=$R_E$, -] (7.0,0) node[ground] {};
  \draw (aux_e) to[short] (8.5,1.2) coordinate (aux_l)
    to[R, l=$R_L$, -] (8.5,0) node[ground] {};

  % Terminal de salida sobre el emisor
  \draw (aux_l) to[short, -o] (9.5,1.2) node[above] {$v_L$} coordinate (out_node);
  \draw (9.5,0) node[ground] {} to[open, v_>=$v_L$] (out_node);
\end{circuitikz}

  \caption{Modelo equivalente de pequeña señal en parámetros híbridos para la etapa en colector común.}
  \label{crkt:pequena_senal}
\end{figure}

A partir de la física de semiconductores, la resistencia dinámica de la juntura emisor-base en directa es:
\begin{equation}
  r_e = h_{ib} = \frac{V_T}{I_{CQ}}
  \label{eq:re_formula}
\end{equation}
donde $V_T = k T / q \approx 25.0\mV$ es el potencial térmico a temperatura ambiente ($T \approx 290\Kelvin$). Evaluando
con la corriente de colector simulada ($I_{CQ} = 1.043\mA$):
\begin{equation}
  r_e = \frac{25.0\mV}{1.043\mA} \approx 23.97\ohm
\end{equation}

La resistencia dinámica de entrada entre base y emisor ($h_{ie}$) resulta:
\begin{equation}
  h_{ie} = (\beta + 1) \cdot r_e = (284 + 1) \cdot 23.97\ohm = 285 \cdot 23.97\ohm \approx 6.83\kohm
  \label{eq:hie_val}
\end{equation}

\section{Cálculo Analítico de Parámetros del Cuadripolo}

\subsection{Impedancia de Entrada (\texorpdfstring{$Z_i$}{Zi})}
La impedancia vista hacia la base del transistor ($R_{i,\mathrm{base}}$) incluye a la resistencia conectada en el emisor por la ganancia de corriente:
\begin{equation}
  R_{i,\mathrm{base}} = h_{ie} + (\beta + 1) \cdot (R_E \parallel R_L)
\end{equation}
\begin{equation}
  R_{i,\mathrm{base}} = 6.83\kohm + 285 \cdot 909.09\ohm = 6.83\kohm + 259.09\kohm \approx 265.92\kohm
\end{equation}

La impedancia de entrada total de la etapa comprende el paralelo con la red de polarización de base ($R_B = R_1 \parallel R_2$):
\begin{equation}
  Z_i = R_B \parallel R_{i,\mathrm{base}} = R_1 \parallel R_2 \parallel \left[ h_{ie} + (\beta + 1)(R_E \parallel R_L) \right]
  \label{eq:Zi_formula}
\end{equation}

Evaluando numéricamente con $R_B = 29.876\kohm$:
\begin{equation}
  Z_i = \frac{29.876\kohm \cdot 265.92\kohm}{29.876\kohm + 265.92\kohm} = \frac{7944.62}{295.80}\kohm \approx 26.86\kohm
  \label{eq:Zi_teorico}
\end{equation}
comprobándose que la etapa ofrece una impedancia de entrada considerablemente mayor a la resistencia del micrófono
($Z_i = 26.86\kohm \ge 2.6 \cdot R_S$), minimizando la atenuación de entrada.

\subsection{Impedancia de Salida (\texorpdfstring{$Z_o$}{Zo})}
La impedancia de salida se evalúa desde los bornes del emisor (desacoplando la carga $R_L$) pasivando la fuente de señal
externa ($v_s = 0$) pero conservando su resistencia interna $R_S$:
\begin{equation}
  Z_o = R_E \parallel \left[ r_e + \frac{R_S \parallel R_B}{\beta + 1} \right]
  \label{eq:Zo_formula}
\end{equation}

Para las condiciones nominales del micrófono ($R_S = 10\kohm$):
\begin{equation}
  R_S \parallel R_B = 10\kohm \parallel 29.876\kohm = \frac{298.76}{39.876}\kohm \approx 7.492\kohm
\end{equation}
\begin{equation}
  Z_o = 1000\ohm \parallel \left[ 23.97\ohm + \frac{7492\ohm}{285} \right] = 1000\ohm \parallel [23.97\ohm + 26.29\ohm] = 1000\ohm \parallel 50.26\ohm \approx 47.85\ohm
  \label{eq:Zo_teorico}
\end{equation}
evidenciando un valor sumamente bajo ($Z_o < 50\ohm \ll R_L = 10\kohm$).

\subsection{Ganancia de Tensión (\texorpdfstring{$A_v$}{Av})}
La tensión alterna desarrollada sobre la carga de salida $R_L$ en bornes del emisor es:
\begin{equation}
  v_L = (\beta + 1) \cdot i_b \cdot (R_E \parallel R_L)
\end{equation}
mientras que la tensión aplicada a la base es:
\begin{equation}
  v_i = i_b \cdot \left[ h_{ie} + (\beta + 1)(R_E \parallel R_L) \right]
\end{equation}

Dividiendo ambas expresiones:
\begin{equation}
  A_v = \frac{v_L}{v_i} = \frac{(\beta + 1)(R_E \parallel R_L)}{h_{ie} + (\beta + 1)(R_E \parallel R_L)} = \frac{R_E \parallel R_L}{r_e + (R_E \parallel R_L)}
  \label{eq:Av_formula}
\end{equation}

Sustituyendo los valores analíticos:
\begin{equation}
  A_v = \frac{909.09\ohm}{23.97\ohm + 909.09\ohm} = \frac{909.09}{933.06} \approx 0.9743
  \label{eq:Av_teorica}
\end{equation}
confirmando una ganancia de tensión no inversora ($0^\circ$) levemente inferior a la unidad debido a la pequeña caída
dinámica en la resistencia intrínseca de emisor $r_e$.

\subsection{Ganancia de Corriente (\texorpdfstring{$A_i$}{Ai})}
La ganancia de corriente representa el cociente entre la corriente entregada a la carga ($i_L$) y la corriente total
absorbida por el amplificador ($i_i$):
\begin{equation}
  i_L = \frac{v_L}{R_L}, \quad i_i = \frac{v_i}{Z_i}
\end{equation}
\begin{equation}
  A_i = \frac{i_L}{i_i} = \frac{v_L / R_L}{v_i / Z_i} = A_v \cdot \frac{Z_i}{R_L}
  \label{eq:Ai_formula}
\end{equation}

Evaluando numéricamente:
\begin{equation}
  A_i = 0.9743 \cdot \frac{26.86\kohm}{10.00\kohm} \approx 2.617
  \label{eq:Ai_teorica}
\end{equation}

\section{Métodos de Medición Experimental en Laboratorio}

Las impedancias dinámicas no pueden relevarse con ohmímetros convencionales debido a que las corrientes continuas
de prueba alterarían la polarización del BJT. Por ello, se implementaron los métodos de inserción resistiva.

\subsection{Medición de \texorpdfstring{$Z_i$, $A_v$ y $A_i$}{Zi, Av y Ai} (Método del Resistor Serie de Entrada)}
En la figura~\ref{crkt:medicion_zi} se muestra el esquema del ensayo experimental.

\begin{figure}[!ht]
  \centering
  \begin{circuitikz}[scale=0.9, transform shape]
  % Generador y resistencia serie
  \draw (0,0) node[ground] {}
    to[sV, l=$v_s$] (0,2.5)
    to[R, l=$R_S$, *-] (2.2,2.5) coordinate (vi_node)
    to[short, -o] (2.2,2.5);

  \draw (0,0) to[short] (2.2,0);
  \draw (0,2.5) to[short, *-o] (0,3.2) node[above] {CH1 ($v_s$)};
  \draw (vi_node) to[short, *-o] (2.2,3.2) node[above] {CH2 ($v_i$)};

  % Bloque amplificador
  \draw (2.7, -0.5) rectangle (5.5, 3.0);
  \node at (4.1, 1.25) [align=center] {Amplificador\\Colector Común};

  \draw (vi_node) to[short] (2.7, 2.5);
  \draw (2.2, 0) to[short] (2.7, 0);

  % Carga y medición de salida
  \draw (5.5, 2.5) to[short, -*] (6.8, 2.5) coordinate (vl_node)
    to[R, l=$R_L$] (6.8, 0) node[ground] {};
  \draw (5.5, 0) to[short] (6.8, 0);
  \draw (vl_node) to[short, -o] (7.5, 2.5) node[right] {$v_L$ (Salida)};
\end{circuitikz}

  \caption{Esquema circuital de medición de $Z_i$, $A_v$ y $A_i$ mediante resistor serie $R_S$ en la entrada.}
  \label{crkt:medicion_zi}
\end{figure}

Se intercala un resistor patrón conocido ($R_S = 27.3\kohm$) entre el generador de funciones y la entrada del circuito.
Midiendo con el osciloscopio la tensión antes de la resistencia ($v_s$) y la tensión en la entrada de la placa ($v_i$),
la corriente de entrada resulta:
\begin{equation}
  i_i = \frac{v_s - v_i}{R_S}
  \label{eq:ii_exp_formula}
\end{equation}

La impedancia de entrada y las ganancias se deducen de forma directa como:
\begin{equation}
  Z_i = \frac{v_i}{i_i} = \frac{v_i \cdot R_S}{v_s - v_i}
  \label{eq:Zi_exp_formula}
\end{equation}
\begin{equation}
  A_v = \frac{v_L}{v_i}
  \label{eq:Av_exp_formula}
\end{equation}
\begin{equation}
  A_i = \frac{i_L}{i_i} = \frac{v_L / R_L}{(v_s - v_i) / R_S}
  \label{eq:Ai_exp_formula}
\end{equation}

\subsection{Medición de \texorpdfstring{$Z_o$}{Zo} (Método de Inyección en Salida)}
En la figura~\ref{crkt:medicion_zo} se ilustra el conexionado de medición de la impedancia de salida.

\begin{figure}[!ht]
  \centering
  \begin{circuitikz}[scale=0.9, transform shape]
  % Bloque amplificador
  \draw (1.2, -0.5) rectangle (4.0, 3.0);
  \node at (2.6, 1.25) [align=center] {Amplificador\\Colector Común};

  % Entrada pasivada (cortocircuitada a masa de CA)
  \draw (0.0, 2.0) to[short] (1.2, 2.0);
  \draw (0.0, 0) to[short] (1.2, 0);
  \draw (0.0, 0) to[short] (0, 2);
  \draw (0.6, 0) node[ground] {};

  % Red de inyección en salida (nodo de emisor)
  \draw (4.0, 2.0) to[short, -*] (5.2, 2.0) coordinate (vo_node)
    to[R, l_=$R_S$] (7.2, 2.0)
    to[sV, l_=$v_s$] (7.2, 0) node[ground] {};
  \draw (4.0, 0) to[short] (7.2, 0);

  \draw (vo_node) to[short, *-o] (5.2, 3.0) node[above] {CH2 ($v_o$)};
  \draw (7.2, 2.0) to[short, *-o] (7.2, 3.0) node[above] {CH1 ($v_s$)};
\end{circuitikz}

  \caption{Esquema de inyección de señal externa para la medición experimental de $Z_o$.}
  \label{crkt:medicion_zo}
\end{figure}

Se cortocircuita la entrada de alterna a masa ($v_i = 0$), se desacopla la carga $R_L$ y se inyecta una señal senoidal
mediante un resistor serie patrón ($R_S = 112\ohm$) sobre el terminal de salida (emisor). La corriente inyectada es:
\begin{equation}
  i_o = \frac{v_s - v_o}{R_S}
\end{equation}
y la impedancia de salida experimental resulta:
\begin{equation}
  Z_o = \frac{v_o}{i_o} = \frac{v_o \cdot R_S}{v_s - v_o}
  \label{eq:Zo_exp_formula}
\end{equation}
